\section*{Capítulo \ref{ch:g7:priorities} --- Prioridad de las operaciones}

\begin{solution}{exo:g7:priorities:1}
$5 + 3 \times 6 = 5 + 18 = 23$.

$(5 + 3) \times 6 = 8 \times 6 = 48$.

$30 - 12 \div 4 = 30 - 3 = 27$.

$(30 - 12) \div 4 = 18 \div 4 = 4.5$.
\end{solution}

\begin{solution}{exo:g7:priorities:2}
$24 - 6 + 2 = 18 + 2 = 20$ (de izquierda a derecha).

$24 \div 6 \times 2 = 4 \times 2 = 8$ (de izquierda a derecha).

$24 - 6 \times 2 = 24 - 12 = 12$ (primero la multiplicación).

$24 \div (6 \times 2) = 24 \div 12 = 2$.
\end{solution}

\begin{solution}{exo:g7:priorities:3}
$7 + 3 \times (10 - 6) = 7 + 3 \times 4 = 7 + 12 = 19$.

$40 - (2 + 3) \times 4 = 40 - 5 \times 4 = 40 - 20 = 20$.

$60 \div \bigl(2 \times (7 - 4)\bigr) = 60 \div (2 \times 3)
= 60 \div 6 = 10$.
\end{solution}

\begin{solution}{exo:g7:priorities:4}
$\dfrac{18 + 6}{8} = \dfrac{24}{8} = 3$; \quad
$\dfrac{30}{7 - 2} = \dfrac{30}{5} = 6$; \quad
$\dfrac{5 \times 8}{4 + 6} = \dfrac{40}{10} = 4$.
\end{solution}

\begin{solution}{exo:g7:priorities:5}
$8 \times 102 = 8 \times (100 + 2) = 800 + 16 = 816$.

$25 \times 99 = 25 \times (100 - 1) = 2500 - 25 = 2475$.

$14 \times 12 + 14 \times 8 = 14 \times (12 + 8) = 14 \times 20 = 280$.

$37 \times 45 - 37 \times 35 = 37 \times (45 - 35) = 37 \times 10 =
370$.
\end{solution}

\begin{solution}{exo:g7:priorities:6}
``El \omterm{ex:g1:subtraction:difference}{diferencia} de $50$ y el \omterm{def:g2:mult:def}{producto} de $6$ por $7$'':
$50 - 6 \times 7 = 50 - 42 = 8$ (no se necesitan corchetes, la prioridad sí
la obra).

``El \omterm{def:g2:mult:def}{producto} del \omterm{def:g1:addition:def}{suma} de $4$ y $5$ por el \omterm{ex:g1:subtraction:difference}{diferencia} de $10$ y
$8$'': $(4 + 5) \times (10 - 8) = 9 \times 2 = 18$.
\end{solution}

\begin{solution}{exo:g7:priorities:7}
$3 \times 2.50 + 2 \times 1.20 = 7.50 + 2.40 = 9.90$ euros.
\end{solution}

\begin{solution}{exo:g7:priorities:8}
(a) sin corchetes: $4 + 2 \times 5 - 1 = 4 + 10 - 1 = 13$.

(b) $(4 + 2) \times 5 - 1 = 30 - 1 = 29$.

(do) $(4 + 2) \times (5 - 1) = 6 \times 4 = 24$.
\end{solution}

\begin{solution}{exo:g7:priorities:9}
$ k = 1$: $12 + 4 \times 1 = 16$ y $16 \times 1 = 16$ --- están de acuerdo
(por eso $ k = 1$ ¡Es una mala prueba!). $ k = 2$:
$12 + 4 \times 2 = 20$ pero $16 \times 2 = 32$: diferente. Entonces los dos
las expresiones son \emph{no} iguales en general. el error
$12 + 4 \times k = 16 \times k $ viene de agregar $12 + 4$ primero,
ignorando la prioridad de $\times $.
\end{solution}

\begin{solution}{exo:g7:priorities:10}
\emph{1.} Precio grupo: $12 + 6 \times n $. Para $ n = 5$:
$12 + 30 = 42$ euros.

\emph{2.} Precio individual: $9n $. Comparar: $ n = 3$: $30$ vs $27$
(individual más barato); $ n = 4$: $36$ vs $36$ (igual); $ n = 5$: $42$ vs
$45$ (grupo más barato). De $5$ personas en, la tarifa de grupo gana
(cada persona extra cuesta $6$ en lugar de $9$, por lo que la brecha no hace más que crecer).
\end{solution}

\begin{solution}{exo:g7:priorities:11}
Posibles respuestas (existen muchas):
\[
24 = 1 \times 2 \times 3 \times 4, \qquad
10 = 1 + 2 + 3 + 4, \qquad
1 = (4 - 3) \times (2 - 1).
\]
\end{solution}

\begin{solution}{pb:g7:priorities:1}
\textbf{1.} $7 \times 102 = 7 \times (100 + 2) = 700 + 14 =
714$; \quad $9 \times 98 = 9 \times (100 - 2) = 900 - 18 = 882$.

\textbf{2.} $17 \times 13 + 17 \times 7 = 17 \times (13 + 7) =
17 \times 20 = 340$; \quad $45 \times 101 - 45 = 45 \times (101
- 1) = 45 \times 100 = 4\,500$.

\textbf{3.} $5 \times 87 = 870 \div 2 = 435$. Legitimidad:
$5 = 10 \div 2$, entonces multiplicando por $5$ se multiplica por $10$
luego dividiendo por $2$ (el orden de $\times $ y $\div $ en el
mismo nivel, de izquierda a derecha, hace el resto).

\textbf{4.} $25 \times 32 = 3\,200 \div 4 = 800$, porque
$25 = 100 \div 4$. Y
$99 \times 47 = (100 - 1) \times 47 = 4\,700 - 47 = 4\,653$:
distributividad sobre el \omterm{ex:g1:subtraction:difference}{diferencia}
(\cref{thm:g7:priorities:distributivity}).

\textbf{5.} $8 \times (5 + 3) = 64$ y $8 \times 5 + 8 \times 3
= 40 + 24 = 64$: coinciden --- eso es distributividad. Pero
$8 \times (5 \times 3) = 8 \times 15 = 120$, mientras
$(8 \times 5) \times (8 \times 3) = 40 \times 24 = 960$: no
partido. El factor se extiende a lo largo del \emph{términos de un \omterm{def:g1:addition:def}{suma}}, no
sobre los factores de un \omterm{def:g2:mult:def}{producto}.

\textbf{6.}
\begin{center}
\renewcommand{\arraystretch}{1.15}
\begin{tabular}{ccl}
$13$ & $24$ & kept ($13$ \omterm{def:g3:numbers:evenodd}{extraño}) \\
$6$ & $48$ & crossed out ($6$ \omterm{def:g3:numbers:evenodd}{incluso}) \\
$3$ & $96$ & kept \\
$1$ & $192$ & kept \\
\end{tabular}
\end{center}
\omterm{def:g1:addition:def}{Suma} de los números correctos supervivientes: $24 + 96 + 192 = 312$, y de hecho $13 \times 24 = 312$.

\textbf{7.} Reducir a la mitad $21$: $21, 10, 5, 2, 1$; duplicación $17$:
$17, 34, 68, 136, 272$. \omterm{def:g3:numbers:evenodd}{Incluso} izquierdas $10$ y $2$ están tachados;
sobrevivientes $17 + 68 + 272 = 357$. Controlar:
$21 \times 17 = 357$.

\textbf{8.} Reducir a la mitad $16$: $16, 8, 4, 2, 1$; duplicación $25$:
$25, 50, 100, 200, 400$. Todos los números restantes excepto el final. $1$
es \omterm{def:g3:numbers:evenodd}{incluso}: sobrevive una sola fila y la respuesta es $400$. una izquierda
columna hecha de duplicaciones puras mantiene sólo su última fila --- la
El método se reduce a duplicar solo.

\textbf{9.} $13 = 12 + 1$, entonces por distributividad
$13 \times 24 = 12 \times 24 + 1 \times 24$; y
$12 \times 24 = 6 \times 48$ desde reducir a la mitad un factor mientras
duplicar el otro deja un \omterm{def:g2:mult:def}{producto} sin cambios
($12 \times 24 = 6 \times 2 \times 24 = 6 \times 48$). El \omterm{def:g3:numbers:evenodd}{extraño}
número izquierdo $13$ dividir una copia del número correcto: el
bancos de primera fila $24$.

\textbf{10.} $6 \times 48 = 3 \times 96$ (mitad-doble, no
banca: $6$ es \omterm{def:g3:numbers:evenodd}{incluso}). $3 \times 96 = 2 \times 96 + 96 =
1 \times 192 + 96$: el \omterm{def:g3:numbers:evenodd}{extraño} $3$ bancos $96$. Finalmente
$1 \times 192 = 192$: el \omterm{def:g3:numbers:evenodd}{extraño} $1$ bancos $192$ y la mesa
termina. Total bancarizado: $24 + 96 + 192 = 312$. Una fila inclina su
número derecho exactamente cuando su número izquierdo es \omterm{def:g3:numbers:evenodd}{extraño} --- entonces el
La receta "mantener las filas impares a la izquierda y agregar" es precisamente la
contabilidad de distributividad.

\textbf{11.} $1 + 4 + 8 = 13$. Para $21$, las filas supervivientes fueron
los de $17$, $68 = 4 \times 17$ y $272 = 16 \times 17$: el
duplicaciones $1 + 4 + 16 = 21$.

\textbf{12.} $25 = 16 + 8 + 1$ y $42 = 32 + 8 + 2$.

\textbf{13.} Tome el mayor número de duplicación que no exceda el
apuntar y restar; lo que queda es cada vez más pequeño
que el número de duplicación que acabamos de usar (de lo contrario, uno mayor
hubiera encajado). Repitiendo, el \omterm{def:g3:division:remainder}{resto} se encoge estrictamente
y debe alcanzar $0$: el proceso se detiene y se utilizan las duplicaciones
son todos diferentes. Para $100$: $64$ encaja ($36$ izquierda), $32$ encaja ($4$ izquierda), $4$ encaja ($0$ izquierda): $100 = 64 + 32 + 4$.

\textbf{14.} tabla de duplicación de $35$: $1 \to 35$, $2 \to 70$,
$4 \to 140$, $8 \to 280$, $16 \to 560$. Desde
$19 = 16 + 2 + 1$, elige esas filas:
$560 + 70 + 35 = 665$. Controlar: $19 \times 35 = 665$.

\textbf{15.} Los mismos tres que el escriba y el campesino:
$24 + 96 + 192$ --- las duplicaciones de $24$ seleccionado por
$13 = 1 + 4 + 8$. Escribir números como \omterm{def:g1:addition:def}{sumas} de duplicaciones es
exactamente escribiéndolos con $0$ arena $1$ s; se cuenta la historia completa
en \cref{pb:g8:powers:1}.
\end{solution}
